\uuid{Tbph}
\exo7id{4422}
\titre{exo7 4422}
\auteur{quercia}
\organisation{exo7}
\datecreate{2010-03-14}
\isIndication{false}
\isCorrection{true}
\chapitre{Série numérique}
\sousChapitre{Autre}
\cpge{mpsi-mp2i.procedes-sommatoires-discrets.a, pcsi.series-numeriques.a, ptsi.series-numeriques.a, mpsi-mp2i.procedes-sommatoires-discrets.b, pcsi.series-numeriques.b, ptsi.series-numeriques.b, mpsi-mp2i.procedes-sommatoires-discrets.c, pcsi.series-numeriques.c, ptsi.series-numeriques.c, mp-mpi.series-numeriques-vectorielles.a, pt.series-numeriques.b, mp-mpi.series-numeriques-vectorielles.b, pc.series-numeriques, psi.series-numeriques.a, pt.series-numeriques.a, bcpst2.series-reelles}
\module{Analyse}
\niveau{L1}
\difficulte{}

\contenu{
\question{
Montrer que la série $\sum_{n=1}^\infty \frac{n^2}{(1+n^2)^2}$ converge.
Calculer une valeur approchée à $10^{-4}$ près de sa somme.
}
\reponse{
$\frac{n^2}{(n^2+1)^2}-\frac{1}{n^2-1}
         = -\frac{3n^2+1}{(n^2+1)^2(n^2-1)} \ge -\frac{4}{n^4}$ pour
         $n\ge 3$.

Donc $S = \sum_{n=1}^N \frac{n^2}{(n^2+1)^2}
        + \sum_{n=N+1}^\infty \frac{1}{n^2-1} + R_N$ avec
     $-\frac{4}{3N^3}\le R_N\le 0$ et
     $\sum_{n=N+1}^\infty \frac{1}{n^2-1} = \frac{N+\frac12}{N(N+1)}$.

Pour $N=25$ on obtient~: $0.76981 < S < 0.76990$.
}
}
